% !TEX program = xelatex
% UBP 论文手稿：Results §3.1–3.5
% 编译：cd manuscript && xelatex results_sec3.tex
% 说明：不依赖 ctexart；使用 article + xeCJK（本机已有）
\documentclass[a4paper,11pt]{article}

\usepackage{geometry}
\geometry{margin=2.2cm}
\usepackage{fontspec}
\usepackage{xeCJK}
\setCJKmainfont{Noto Sans CJK SC}
\setmainfont{DejaVu Serif}
\usepackage{graphicx}
\usepackage{amsmath,amssymb}
\usepackage{xcolor}
\usepackage{hyperref}
\usepackage{caption}
\usepackage{setspace}
\onehalfspacing

\captionsetup{font=small,labelfont=bf,skip=6pt}
\definecolor{linkblue}{RGB}{0,51,128}
\hypersetup{colorlinks=true,linkcolor=linkblue,urlcolor=linkblue}

\title{结果（摘录）}
\author{}
\date{}

\begin{document}
\maketitle

\noindent{\footnotesize
除非另有说明，以下结果基于 THINGS-EEG 的独立测试集
（$n=10$ 被试；每位被试 200 个测试图像；每个图像的 GT 为 80 次重复试次平均）。
值得注意的是，该 200 个测试概念均未出现在 1654 个训练类别中
（训练--测试类别交集为空），因此本节评估对应类别层面的零样本泛化，
而非对已见类别新图像的插值。
第~\ref{sec:alljoined}~节另报告 Alljoined-1.6M 的跨数据集评估
（采用相同的未见类别测试划分）。
预测脑电与真实脑电分别记为 Pred 与 GT。
正文中的汇总值默认为跨被试均值$\pm$标准差；
图中误差线或阴影所表示的 SD/SEM 均在相应图注中单独说明。
}

%======================================================================
\setcounter{section}{2}
\section{结果}
\label{sec:results}

%----------------------------------------------------------------------
\subsection{预测脑电保留视觉诱发响应的时空组织}
\label{sec:spatiotemporal}

Pred 在未见训练类别的独立测试图像上恢复了稳定的视觉诱发响应，并保留其时空组织。
将 200 个零样本测试图像与 10 名被试汇总后，Pred 在连续 $100\,\mathrm{ms}$ 时间窗内
复现了 GT 随时间变化的头皮极性及其后部优势
（Fig.~\ref{fig:spatiotemporal}a）。
补充分析从三个层次支撑同一结论：逐时间点相关显示对应关系贯穿刺激后时程，
通道$\times$时间热图保留完整 ERP 条带与极性转换，
被试级汇总指标则给出幅值误差与波形保真的定量读数
（Supplementary Fig.~\ref{fig:sfig1}a,b,e）。
全局 RMSE 为 $0.179\pm0.020$，pooled Pearson 相关为 $r=0.525\pm0.078$，
mean channel $r$ 为 $0.299\pm0.030$
（Supplementary Fig.~\ref{fig:sfig1}e）；
由同一误差矩阵得到的全局 MAE 为 $0.135\pm0.013$。
因此，Pred 与 GT 的一致不仅体现为较低幅值误差，也体现为整体波形结构的对应。

该对应并非由单一瞬时峰值驱动。
逐时间点 Pred--GT 相关在 $0$--$996\,\mathrm{ms}$ 内始终为正，
跨时间均值为 $r=0.506$，并在 $108\,\mathrm{ms}$ 达到描述性峰值
（$r=0.659$；Supplementary Fig.~\ref{fig:sfig1}a）。
这表明模型保留的是贯穿刺激后时程的响应结构。

预测保真度在头皮上并非均匀分布，而表现为前部较低、后部较高的空间梯度。
63 个通道总体均值为 $r=0.299$，枕区 8 个通道达到 $r=0.749$；
Oz 最高（$r=0.907$），F1 最低（$r=0.013$）
（Fig.~\ref{fig:spatiotemporal}b）。
在互斥的五区映射下，Frontal、Central、Temporal、Parietal 和 Occipital
的描述性均值分别为 $0.164$、$0.160$、$0.254$、$0.494$ 和 $0.749$
（Fig.~\ref{fig:spatiotemporal}b,c）。
由于未进行脑区间推断检验，该结果仅解释为描述性空间梯度。

这一后部优势进一步落实到可辨识的 ERP 波形。
在全通道平均相关最高的示例被试 sub-08 中，
200 个测试图像平均后的 Oz、PO4、O1、POz、PO8 和 PO3 波形
在峰谷时相和相对幅度上保持一致，通道相关为 $0.79$--$0.90$
（Fig.~\ref{fig:spatiotemporal}d）；该面板仅作代表性展示，不用于被试间推断。
同一被试的单图像示例进一步表明，该对应并不依赖跨图像平均：
5 张固定测试图像的全通道平均相关为 $r=0.47$--$0.51$，
所示后部通道波形相关为 $0.87$--$0.94$，
对应 log-PSD 相关为 $0.95$--$0.98$
（Supplementary Fig.~\ref{fig:sfig2}）。
这些单图像结果仍以每张图像 80 次重复试次平均的 GT 为参照，
因此说明的是图像层面稳定响应的可预测性，而非单试次预测。
在群体层面，枕区 grand-average ERP 的 Pred--GT 波形相关为 $r=0.972$，
被试内相关再平均为 $0.968\pm0.020$
（Fig.~\ref{fig:spatiotemporal}e）。

枕区波形一致性延伸至早期视觉 ERP 成分。
C1（$50$--$90\,\mathrm{ms}$）、P1（$80$--$130\,\mathrm{ms}$）和
N1（$140$--$200\,\mathrm{ms}$）的波形相关分别为
$0.723\pm0.130$、$0.726\pm0.157$ 和 $0.615\pm0.093$；
跨图像幅度相关为 $0.310\pm0.113$、$0.531\pm0.144$ 和 $0.505\pm0.087$，
峰值潜伏期 MAE 为 $5.7\pm2.2$、$7.0\pm3.2$ 和 $13.2\pm5.4\,\mathrm{ms}$
（Fig.~\ref{fig:spatiotemporal}f）。
综上，Pred 保留了从全局时程、头皮空间梯度到毫秒尺度早期成分的多层时空组织；
上述评估均针对试次平均响应，而非单试次 EEG 变异。

\begin{figure}[!t]
  \centering
  \includegraphics[width=\textwidth]{figures/fig1.png}
  \caption{\textbf{预测脑电与真实脑电的时空一致性。}
  \textbf{a}, GT 与 Pred 的 grand-average ERP 地形图（$100\,\mathrm{ms}$ 时间窗）。
  \textbf{b}, 按五个互斥头皮区域分组的 63 通道 Pred--GT Pearson $r$
  （柱：跨被试均值；误差线：跨被试 SD）；红虚线表示通道均值。
  \textbf{c}, 使用与 \textbf{b} 相同区域映射得到的逐通道预测质量分布；
  散点中不显示超出各区域 $1.5\times\mathrm{IQR}$ 的 Tukey 离群值，
  均值虚线及正文统计仍基于完整数据。该面板用于描述空间梯度，
  未进行脑区间推断检验。
  \textbf{d}, 全通道平均相关最高的示例被试（sub-08）由
  200 个测试图像平均得到的后部通道波形（GT 为实线，Pred 为虚线）。
  \textbf{e}, 枕区 grand-average ERP；曲线和阴影表示跨被试 mean$\pm$SEM。
  \textbf{f}, 枕区 C1、P1 和 N1 成分的波形相关、跨图像幅度相关及峰值潜伏期 MAE
  （符号和误差线：跨被试 mean$\pm$SEM；浅色点：被试）。
  除 \textbf{d} 外，$n=10$ 独立被试。}
  \label{fig:spatiotemporal}
\end{figure}

%----------------------------------------------------------------------
\subsection{预测脑电保留频谱形状及频带特异的空间结构}
\label{sec:frequency}

Pred 不仅恢复了时域波形，也恢复了频谱形状及其后部空间组织。
被试层面全脑 $0.5$--$45\,\mathrm{Hz}$ log-PSD 相关为 $0.968\pm0.015$，
五频带 log 带功率相关均值为 $0.685\pm0.037$
（Supplementary Fig.~\ref{fig:sfig1}e）。
跨被试平均后的全脑 PSD 曲线相关为 $r=0.973$，
且 Pred 与 GT 均在 $9.77\,\mathrm{Hz}$ 出现 $\alpha$ 峰
（Supplementary Fig.~\ref{fig:sfig1}c）；
枕区组平均 log-PSD 相关进一步达到 $r=0.992$，峰值位置不变
（Fig.~\ref{fig:frequency}d）。
Pred 在高频端功率偏低，但因未对频段功率差或峰频做推断检验，
这里仅将其视为描述性偏差。

平均谱包络的一致性伴随着频带功率空间分布的保留。
在全频段及 $\delta$、$\theta$、$\alpha$、$\beta$、$\gamma$ 频带中，
Pred 与 GT 呈现对应的 log 带功率拓扑，
其中 $\theta$/$\alpha$ 的后部优势最为明显
（Fig.~\ref{fig:frequency}a）。
在各频带内，先于被试内对图像$\times$通道的 log 带功率计算 Pearson 相关，
再跨被试平均，得到
$\delta$：$0.366\pm0.068$，$\theta$：$0.855\pm0.024$，
$\alpha$：$0.800\pm0.057$，$\beta$：$0.803\pm0.048$，
$\gamma$：$0.601\pm0.125$。
因此，频率结构的保留同时体现在平均谱形与图像/通道间的带功率变异上。

频带—通道相似矩阵进一步将这一保真度定位到逐通道波形。
按 Frontal、Central、Temporal、Parietal 和 Occipital 分组后，
五个频带、63 通道的 Pred--GT 波形相关总体均值为 $r=0.329$
（Fig.~\ref{fig:frequency}c）。
$\theta$ 的全脑与枕区均值最高（$0.514$ 和 $0.860$），
其次为 $\alpha$（$0.375$ 和 $0.792$）及 $\delta$（$0.341$ 和 $0.604$），
$\gamma$ 最低（$0.128$ 和 $0.376$）；最大值出现在 $\theta$@Oz（$r=0.968$）。
该按脑区分组的分布延续了第~\ref{sec:spatiotemporal}~节的后部梯度，
表明较稳定保留的频率结构主要集中于顶枕通道的中低频成分。

高相关仍可能来自共享的视觉诱发模板，而非图像特异信息。
为此，我们计算 True--Shuffle band fidelity：
将 Pred 与对应 GT 的频带限制波形相关定义为 True，
并与 50 次随机错配 GT 的平均相关（Shuffle）比较；
先在每名被试内对 200 个图像平均，再以单侧 Wilcoxon 符号秩检验评估
相对增益 $(\mathrm{True}-\mathrm{Shuffle})/\mathrm{True}$ 是否超过 $3\%$ 下限。
在全脑尺度，All、$\delta$、$\alpha$ 和 $\beta$ 均超过该下限
（Fig.~\ref{fig:frequency}b）：
All 为 $0.611\pm0.018$ vs.\ $0.562\pm0.017$，
$\delta$ 为 $0.421\pm0.024$ vs.\ $0.309\pm0.026$
（两者均 $W=55$，$P=9.8\times10^{-4}$，$r_{\mathrm{rb}}=1.00$；均值$\pm$SEM）；
$\alpha$（$W=54$，$P=0.0020$，$r_{\mathrm{rb}}=0.964$）和
$\beta$（$W=52$，$P=0.0049$，$r_{\mathrm{rb}}=0.891$）方向一致。
四项检验均通过六频带 Bonferroni 阈值 $\alpha=0.05/6=0.0083$。
相比之下，$\theta$ 虽绝对相关较高，但 True 与 Shuffle 同时偏高，
其 $3.4\%$ 相对增益未超过门槛（$W=39$，$P=0.14$）；
$\gamma$ 同样未达标（$W=0$，$P=1.00$）。

枕区 True--Shuffle band fidelity 复现了同一选择性，且整体相关更高
（Supplementary Fig.~\ref{fig:sfig1}d）。
All（$0.859\pm0.010$ vs.\ $0.795\pm0.013$）和
$\delta$（$0.624\pm0.014$ vs.\ $0.383\pm0.025$）
均超过 $3\%$ 下限（均 $W=55$，$P=9.8\times10^{-4}$，$r_{\mathrm{rb}}=1.00$），
$\alpha$ 与 $\beta$ 亦达标（均 $W=54$，$P=0.0020$，$r_{\mathrm{rb}}=0.964$）；
$\theta$（$W=30$，$P=0.42$）和 $\gamma$（$W=16$，$P=0.88$）未达标。
达到阈值的四项均通过 Bonferroni 校正。
因此，Pred 不仅恢复了平均谱包络及其后部空间组织，
还在 All、$\delta$、$\alpha$ 和 $\beta$ 频段保留了超出随机错配基线的图像匹配信息；
该控制量化的是正确配对相对随机错配的增益，
并不能完全分离所有测试图像共享的视觉诱发成分。

\begin{figure}[!t]
  \centering
  \includegraphics[width=\textwidth]{figures/fig2.png}
  \caption{\textbf{预测脑电的频域保真性。}
  \textbf{a}, GT 与 Pred 在全频段及五个经典频带上的 log 带功率地形。
  \textbf{b}, 全脑 True--Shuffle band fidelity：
  匹配（True）与随机错配（Shuffle）的频带限制波形相关。
  柱和误差线表示跨被试 mean$\pm$SEM，散点表示被试；
  星号按相对增益超过 $3\%$ 下限的单侧 Wilcoxon $p$ 标注（未校正）；
  正文另报告六频带 Bonferroni 校正后的推断结论。
  \textbf{c}, 五个频带、63 个通道的 Pred--GT 波形相关；
  横轴通道按 Frontal、Central、Temporal、Parietal 和 Occipital 分组排列。
  \textbf{d}, 枕区平均 PSD；GT 与 Pred 均在 $9.77\,\mathrm{Hz}$ 达到
  $\alpha$ 峰。$n=10$ 独立被试。}
  \label{fig:frequency}
\end{figure}

%----------------------------------------------------------------------
\subsection{预测脑电保留刺激相关的表征对应与几何结构}
\label{sec:representational}

时空与频谱一致仍可能主要反映共享的视觉诱发成分；
Pred 是否保留区分具体刺激的高维表征，需要在预定嵌入空间中直接检验。
我们将每名被试的 GT 与 Pred 分别输入同一冻结的被试内脑编码器，
并在跨模态图像对应、同刺激 GT--Pred 配对和概念间表征几何三个层面加以评估。
分析仍使用 200 个未见训练类别的测试概念试次平均 EEG；
编码器输出为归一化的 1024 维嵌入，并与 CLIP 图像特征处于训练时建立的共享空间。
作为总览，Pred EEG--CLIP semantic cosine 为 $0.323\pm0.011$，
被试内 GT--Pred RDM RSA 为 $\rho=0.469\pm0.029$
（Supplementary Fig.~\ref{fig:sfig1}e）。
这些汇总值尚不能区分对应来自对角刺激配对、同图像 GT--Pred 接近，
还是概念间几何结构的保留，因此以下分别检验。

Pred 在共享嵌入空间中保留了明确的刺激对应。
跨 10 名被试逐元素平均后，GT 与 Pred 的 EEG--图像余弦相似度矩阵
均沿正确配对对角线增强，并在按 animal、food、vehicle、tool 和 others
排序后保留可见的类内块结构（Fig.~\ref{fig:representational}a）。
组平均矩阵中，GT 的对角与非对角相似度均值分别为 $0.176$ 和 $0.053$，
Pred 为 $0.323$ 和 $0.054$。
因此，Pred 的对角增强并非对所有图像普遍抬高相似度。
不过，由于生成器的语义训练目标与此处评价均依赖 CLIP 对齐空间，
Pred 较高的对角相似度不能解释为比真实 EEG 具有更强的内在语义表征，
而应视为对模型目标空间内跨模态对应的检验。

同刺激 GT--Pred 嵌入配对进一步表明，该对应包含图像特异信息。
在 2,000 个同被试、同图像配对中，余弦相似度均值为 $0.414$
（配对层面 SD $=0.125$；被试级均值 $0.414\pm0.037$），
而 4,980 个被试内随机错配均值仅为 $0.043$，均值差 $\Delta\mu=0.372$
（Fig.~\ref{fig:representational}b）。
这一描述性分离提示 Pred 与 GT 的接近不仅来自被试内共享的平均嵌入结构，
还包含与具体测试图像绑定的信息；
由于随机对照为抽样非配对分布且未做被试层面推断检验，这里不表述为显著差异。

该刺激对应还扩展到概念间表征几何。
以 GT 与 Pred 编码器嵌入分别构建 $200\times200$ RDM
（元素为 $1-\cos$）后，跨被试平均的两张 RDM 呈现相似的概念间距离结构
（Fig.~\ref{fig:representational}c），
其 19,900 个非冗余上三角元素的 Spearman 相关为 $\rho=0.691$
（Fig.~\ref{fig:representational}d）。
需要区分的是，该值由两张组平均 RDM 计算；
先在每名被试内计算 RSA 再汇总所得为 $\rho=0.469\pm0.029$。

置换检验表明，上述被试级 RSA 不能由任意图像配对解释。
打乱每名被试的图像--Pred 对应并重建 Pred RDM 后，
在 1,000 次单侧置换中，观测到的被试级 RSA 均值（$0.469\pm0.029$）
高于组水平零分布（$0.00004\pm0.00250$，$P=0.001$），
且 10/10 名被试均达到置换分辨率下的 $P=0.001$。
作为补充诊断，2,000 个 GT 与 2,000 个 Pred 嵌入的联合 UMAP
密度重叠为 $0.82$（Supplementary Fig.~\ref{fig:sfig3}b）；
Pred--GT RDM 残差上三角均值为 $0.061$（SD $=0.100$），
提示 Pred 中概念间距离存在轻度描述性扩张，但主要距离排序仍得以保留
（Supplementary Fig.~\ref{fig:sfig3}c）。
五类检索案例提供直观示例（Supplementary Fig.~\ref{fig:sfig3}a），
但其来自单一被试且经 GT 与 Pred 均进入 Top-5 的预筛选，
故不作为群体检索性能的推断证据。
综上，Pred 在预定脑编码器空间中保留了刺激特异的 GT--Pred 对应及概念间几何；
置换 RSA 构成这一几何对应超出随机图像配对的主要统计证据。

\begin{figure}[!t]
  \centering
  \includegraphics[width=\textwidth]{figures/fig3.png}
  \caption{\textbf{预测脑电的刺激特异表征对应与几何结构。}
  \textbf{a}, GT 和 Pred 脑编码器嵌入与 200 个 CLIP 图像特征之间的
  余弦相似度矩阵；矩阵先在每名被试内计算，再跨 10 名被试逐元素平均。
  行和列按五类测试概念排序，两张矩阵使用同一色标。
  \textbf{b}, 同被试、同图像的 GT--Pred 嵌入配对
  （Matched；$n=2{,}000$）与被试内随机错配
  （Random；$n=4{,}980$）的余弦相似度密度；
  虚线表示各分布均值，$\Delta\mu$ 表示均值差。该面板为描述性比较。
  \textbf{c}, 由 GT 和 Pred 脑编码器嵌入构建并跨被试平均的
  $200\times200$ RDM；不相似度定义为 $1-\cos$。
  \textbf{d}, 两张组平均 RDM 的 19,900 个非冗余上三角元素对应关系；
  颜色表示点密度，虚线为单位线，标注值为 Spearman $\rho$。
  全部分析基于 $n=10$ 独立被试和每名被试 200 个测试概念。}
  \label{fig:representational}
\end{figure}

%----------------------------------------------------------------------
\subsection{跨数据集评估复现时空、频谱与表征对齐}
\label{sec:alljoined}

同一套时空、频谱与表征对应也能在消费级 EEG 上同向复现。
我们在 Alljoined-1.6M 上复用与 THINGS-EEG 相同的评价流程
（32 通道 Emotiv Flex 2 蒙太奇；相同 THINGS 图像刺激与
未见训练类别的 200 概念零样本测试划分）。
由于消费级记录信噪比更低、被试间变异更大，
跨数据集分析并非直接纳入全部 20 名被试，
而是先依据冻结脑编码器的刺激可分性进行预筛选。
具体地，我们在全部 20 名被试上训练被试内脑编码器后，
以测试集 200 类图像检索的 Top-5 准确率 $\ge 20\%$ 作为预定义阈值。
该阈值等于随机基线（$5/200=2.5\%$）的 8 倍，
并落在被试表现的自然分界处：
超过阈值的 5 名被试 Top-5 为 $20.5\%$--$44.0\%$，
其余 15 名均 $\le 16.0\%$。
据此筛选出的 sub-06、sub-12、sub-13、sub-14 和 sub-18
构成高可分性子集；后续生成器训练与评价仅在该子集上进行。
阈值在生成器评估之前设定，目的是保证语义读出
（semantic cosine、RSA）所依赖的脑编码器--CLIP 空间本身具有可检测的刺激对应，
而非依据生成器表现事后挑样。

在这些预筛选被试上，Pred 保留了可识别的时空与频谱结构。
被试级汇总显示，RMSE 为 $0.356\pm0.069$，
pooled $r$ 为 $0.257\pm0.071$，mean channel $r$ 为 $0.208\pm0.063$，
全脑 log-PSD 相关为 $0.944\pm0.013$，五频带 bandpower 相关均值为 $0.472\pm0.052$
（Supplementary Fig.~\ref{fig:sfig4}a）。
通道$\times$时间热图与枕区 ERP 表明 Pred 复现了主要振荡组织
（Supplementary Fig.~\ref{fig:sfig4}b,d）；
枕区 grand-average 波形相关为 $r=0.984$。
按通道名称划分的四区汇总中，Frontal、Central、Parietal 和 Occipital
的描述性均值分别为 $0.094$、$0.128$、$0.186$ 和 $0.418$，
后部优势与 THINGS-EEG 方向一致
（Supplementary Fig.~\ref{fig:sfig4}c）。

表征层面同样保留刺激特异对应。
Pred EEG--CLIP semantic cosine 为 $0.235\pm0.010$，
被试内 GT--Pred RDM RSA 为 $\rho=0.286\pm0.070$
（Supplementary Fig.~\ref{fig:sfig4}a）。
在 1,000 个同被试、同图像配对中，匹配嵌入余弦均值高于随机错配
（$0.239$ vs.\ $0.058$，$\Delta\mu=0.181$；
Supplementary Fig.~\ref{fig:sfig4}f）；
组平均 RDM 的 Spearman 相关为 $\rho=0.524$
（Supplementary Fig.~\ref{fig:sfig4}g）。
因此，尽管消费级记录的绝对波形指标低于研究级 THINGS-EEG，
时空组织、频谱形状与刺激相关表征对齐仍在外部数据集上同向成立。
该结论限于预筛选的高可分性子集，
不能外推为 Alljoined 全部 20 名被试的群体表现。

%----------------------------------------------------------------------
\subsection{结构与损失消融揭示预测质量的不同约束来源}
\label{sec:ablation}

在 THINGS-EEG 上建立多层保真、并在 Alljoined 上观察到同向复现之后，
我们进一步通过结构和损失消融分析这些能力分别来自哪些设计约束，
并检验单一训练目标是否足以同时优化全部读出。
我们关注三个互补指标：mean channel Pearson $r$（时域波形）、
bandpower correlation（频带功率结构）和
semantic cosine（Pred EEG 在脑编码器--CLIP 空间中与对应图像的跨模态对齐）。
除非另有说明，本节正文中的 $\pm$ 为跨被试 SD；Fig.~\ref{fig:ablation} 误差线为 SEM。

视觉输入形式主要约束语义对齐，而非波形或频带指标。
以 sharp CLIP 替代 foveated CLIP 后，semantic cosine 从
$0.323\pm0.011$ 降至 $0.248\pm0.010$
（Fig.~\ref{fig:ablation}a），
该下降在扩展消融的被试配对 Wilcoxon 检验中显著
（$P=9.8\times10^{-4}$；Bonferroni $\alpha=0.05/45=0.0011$）。
相比之下，Pearson $r$（$0.309\pm0.036$）和 bandpower correlation（$0.691\pm0.035$）
相对最终模型（$0.299\pm0.030$ 与 $0.685\pm0.037$）并未下降。
因此，foveated 表征的主要贡献是提供更适合共享嵌入空间的语义约束。

时间生成结构则更直接影响波形与频谱保真。
去除 dilated TCN 后，Pearson $r$ 降至 $0.253\pm0.034$，
bandpower correlation 降至 $0.646\pm0.037$；
其中 Pearson $r$ 的下降在结构消融中显著
（$P=9.8\times10^{-4}$；Bonferroni $\alpha=0.05/20=0.0025$），
而 bandpower 下降未通过该阈值（$P=0.0029$）。
完全移除 Transformer 后下降更广泛：
Pearson $r$ 为 $0.197\pm0.026$，bandpower correlation 为 $0.610\pm0.056$，
两项均显著（$P=9.8\times10^{-4}$ 和 $P=0.0020$；Fig.~\ref{fig:ablation}a）。
这表明 dilated TCN 主要支撑局部与多尺度时程建模，
而 Transformer 对较长程时频结构整合更为关键。
容量或注意力变体并不呈单调规律：
例如 hidden dimension 为 128 时 Pearson $r$ 较高（$0.324\pm0.029$），
但 bandpower correlation 降至 $0.655\pm0.019$；
hidden dimension 为 512 时波形与频带指标均下降。
因此，不同读出并不由模型容量单调决定。

损失组消融进一步揭示优化目标之间的相互制约。
在单阶段训练中，去除 time loss 几乎消除时域波形对应
（Pearson $r=0.036\pm0.021$），semantic cosine 反而保持较高
（$0.363\pm0.008$；Fig.~\ref{fig:ablation}b）。
相反，去除 semantic loss 后，semantic cosine 降至 $0.068\pm0.011$，
而 Pearson $r$（$0.262\pm0.025$）和 bandpower correlation（$0.594\pm0.075$）
并未发生同等幅度下降。
因此，时间损失与语义损失分别约束其对应读出，
但单独强化其一并不会同步改善其他维度。

频域目标呈现不同模式。
相对单阶段 Full，去除 frequency loss 并未造成清晰的 bandpower correlation 下降
（$0.580\pm0.025$ vs.\ $0.593\pm0.072$），
同时 Pearson $r$ 升至 $0.313\pm0.034$。
这说明在单阶段联合优化中，频域项易被时域与语义目标抵消，
也解释了最终模型采用两阶段训练的原因：
先建立稳定的时间与语义映射，再通过 spectral refinement 强化频带结构。
与单阶段 Full 相比，两阶段训练将 bandpower correlation 从
$0.593\pm0.072$ 提升至 $0.685\pm0.037$，
Pearson $r$ 从 $0.251\pm0.025$ 提升至 $0.299\pm0.030$，
而 semantic cosine 从 $0.345\pm0.008$ 小幅回落至 $0.323\pm0.011$。
这种非一致变化表明时域、频域与语义目标之间存在实际权衡：
最终模型并非在每一单项上取最大值，
而是在保留语义对齐的同时提高波形与频域保真，形成前述多层结果的折中解。

上述多层保真并不依赖大容量生成模型。
最终 EEG 生成器仅含 $9.16\times10^{6}$ 个可训练参数（FP32 约 $35\,\mathrm{MiB}$）；
推理时另需冻结的 CLIP RN50 视觉编码器（约 $3.83\times10^{7}$ 参数），
二者合计约 $4.75\times10^{7}$ 参数。
在消费级 GPU（NVIDIA GeForce RTX 5060 Ti）上，
对单张测试图像执行完整链路——foveated 模糊预处理（核大小 $51$）、
CLIP 图像编码与 EEG 生成——平均耗时 $8.4\,\mathrm{ms}$
（中位数 $7.5\,\mathrm{ms}$；$n=50$ 次重复，不含磁盘读图，batch size $=1$）。
若图像特征已预先缓存，仅生成器前向约为 $2.4\,\mathrm{ms}$。
因此，第~\ref{sec:spatiotemporal}--\ref{sec:alljoined}~节所报告的
时空、频谱与表征对应，是在紧凑参数规模与毫秒级延迟下取得的。

\begin{figure}[!t]
  \centering
  \includegraphics[width=\textwidth]{figures/fig4.png}
  \caption{\textbf{生成器结构和损失目标的消融分析。}
  \textbf{a}, 预测器结构消融。比较最终模型（Ours）与去除或改变关键结构后的
  mean channel Pearson $r$、bandpower correlation 和 semantic cosine。
  w/o Fovea 表示以 sharp CLIP 输入替代 foveated CLIP 输入；
  w/o Dilated 表示去除 dilated TCN；
  w/o TF 表示移除 Transformer；
  w/o attn 表示去除 self-attention；
  H=128 和 H=512 表示改变隐藏维度。
  \textbf{b}, 损失组消融。Ours 为两阶段训练后的最终模型；
  Single-stage 为同一组时间、频域和语义损失的单阶段联合训练；
  w/o time、w/o freq 和 w/o semantic 分别移除时间、频域或语义损失组。
  柱和误差线表示跨 10 名被试的 mean$\pm$SEM；
  蓝色虚线表示最终模型的对应指标。}
  \label{fig:ablation}
\end{figure}

\begingroup
\renewcommand{\thefigure}{S1}
\begin{figure}[!p]
  \centering
  \includegraphics[width=0.92\textwidth]{figures/Sfig1.png}
  \caption{\textbf{预测质量的补充评估。}
  \textbf{a}, 逐时间点 Pred--GT Pearson $r$；
  实线和阴影分别表示跨被试均值和 SD。
  \textbf{b}, GT 与 Pred 的 grand-average ERP 热图，
  展示通道$\times$时间结构及其极性转换。
  \textbf{c}, 全脑平均 PSD，用于检验平均谱形是否被保留。
  \textbf{d}, 枕区 True--Shuffle band fidelity（occ.）：
  匹配（True）与随机错配（Shuffle）的频带限制波形相关，
  作为 Fig.~\ref{fig:frequency}b 的枕区对照，
  用于评估正确图像配对相对于随机错配的频带特异增益。
  \textbf{e}, 被试级汇总指标：全图像$\times$通道$\times$时间 RMSE；
  pooled $r$（相同维度联合展平）；mean channel $r$
  （每通道在图像$\times$时间上池化后对 63 通道平均）；
  PSD $r$（全脑 $0.5$--$45\,\mathrm{Hz}$ log-PSD）；
  Bandpower $r$（五频带 log 带功率相关的均值）；
  Cos（Pred EEG--CLIP semantic cosine）；
  RSA（被试内 GT--Pred RDM Spearman $\rho$）。
  \textbf{d}、\textbf{e} 中柱和误差线表示跨被试 mean$\pm$SEM，散点表示被试。
  \textbf{a}--\textbf{b} 展示时空细节；
  \textbf{c}--\textbf{d} 展示频域结果；
  \textbf{e} 汇总误差、波形、频谱和表征指标。$n=10$ 独立被试。}
  \label{fig:sfig1}
\end{figure}
\endgroup

\begingroup
\renewcommand{\thefigure}{S2}
\begin{figure}[!p]
  \centering
  \includegraphics[width=0.92\textwidth]{figures/Sfig2.png}
  \caption{\textbf{单图像层面的波形和频谱预测示例。}
  示例被试 sub-08 中 5 张固定测试图像的单图像预测结果。
  左侧显示测试图像、概念名称和该图像在 63 通道上的平均 Pearson $r$。
  中间三列显示每张图像波形相关最高的 3 个后部通道；
  蓝色实线为 GT，红色虚线为 Pred，标题中标注对应通道的波形相关。
  右侧显示该图像 log-PSD 相关最高通道的频谱曲线，
  标题中标注对应的 PSD 相关。该图仅用于说明图像层面的代表性预测，
  不用于被试间推断。}
  \label{fig:sfig2}
\end{figure}
\endgroup

\begingroup
\renewcommand{\thefigure}{S3}
\begin{figure}[!p]
  \centering
  \includegraphics[width=0.92\textwidth]{figures/Sfig3.png}
  \caption{\textbf{表征对应的补充诊断。}
  \textbf{a}, 示例被试 sub-01 的五类图像检索案例。
  每行以左侧图像对应的 GT 或 Pred EEG 嵌入为查询，
  在 200 个 CLIP 图像特征中显示余弦相似度最高的五个结果。
  每类案例从 GT 与 Pred 均将正确图像检索至 Top-5 的候选中
  以固定随机种子抽取，因此该面板仅作定性展示。
  \textbf{b}, 将 10 名被试的 2,000 个 GT 与 2,000 个 Pred 脑编码器嵌入
  拼接后，经标准化、50 维 PCA 和同一次 UMAP 拟合得到的联合二维投影；
  实线和虚线分别包围 GT 与 Pred 的 $80\%$ KDE 概率质量，
  KDE overlap 为两归一化密度的重叠系数。
  \textbf{c}, 跨被试平均的 Pred RDM 减去 GT RDM；
  色标表示 $\Delta(1-\cos)$，并对称设为非冗余残差绝对值的
  $99.5$ 百分位（超出范围的值在显示时饱和）。
  \textbf{b}、\textbf{c} 为描述性降维和残差诊断，不承担独立推断。}
  \label{fig:sfig3}
\end{figure}
\endgroup

\begingroup
\renewcommand{\thefigure}{S4}
\begin{figure}[!p]
  \centering
  \includegraphics[width=0.92\textwidth]{figures/Sfig4.png}
  \caption{\textbf{Alljoined 消费级 EEG 上的跨数据集评估。}
  分析基于 Alljoined-1.6M 中编码器测试集 Top-5 $\ge 20\%$
  （等于 8 倍随机基线）的 5 名预筛选被试：
  sub-06、sub-12、sub-13、sub-14、sub-18；
  每位被试 200 个未见训练类别的测试图像，GT 为 80 次重复试次平均。
  \textbf{a}, 被试级汇总指标：RMSE、pooled $r$、mean channel $r$、
  全脑 PSD $r$、Bandpower $r$、Pred EEG--CLIP semantic cosine 与
  被试内 GT--Pred RDM RSA；柱和误差线为 mean$\pm$SEM，散点为被试。
  \textbf{b}, GT 与 Pred 的 grand-average ERP 热图。
  \textbf{c}, 按 Frontal、Central、Parietal 和 Occipital 汇总的逐通道波形相关；
  虚线为总体均值。
  \textbf{d}, 枕区 grand-average ERP（曲线与阴影：跨被试 mean$\pm$SEM）。
  \textbf{e}, 枕区平均 PSD。
  \textbf{f}, 同被试同图像 GT--Pred 嵌入配对与随机错配的余弦相似度密度；
  $\Delta\mu$ 为均值差。
  \textbf{g}, 组平均 GT 与 Pred RDM 上三角元素对应关系，标注 Spearman $\rho$。
  该图用于外部验证，不与 THINGS-EEG 主结果做绝对数值优劣比较；
  结论限于上述预筛选子集。}
  \label{fig:sfig4}
\end{figure}
\endgroup

\end{document}
